\documentclass{article}
\usepackage{graphicx} % Required for inserting images
\usepackage{xcolor}

\usepackage{amsmath}
\usepackage{amsthm}
\usepackage{thmtools} 
\usepackage{cleveref}


% Number equations within sections (instead of chapters)
\numberwithin{equation}{section}

% Declare theorem environments anchored to sections
\declaretheorem[name=Theorem,numberwithin=section]{theorem}
\declaretheorem[name=Lemma,numberlike=theorem]{lemma}
\declaretheorem[name=Proposition,numberlike=theorem]{proposition}
\declaretheorem[name=Corollary,numberlike=theorem]{corollary}
\declaretheorem[name=Definition,numberlike=theorem]{definition}
\declaretheorem[name=Assumption,numberlike=theorem]{assumption}
\declaretheorem[name=Example,numberlike=theorem]{example}
\declaretheorem[name=Remark,numberlike=theorem]{remark}

% For restatable theorems 
\declaretheorem[name=Theorem,numberwithin=section]{restatabletheorem}

\title{Cooperative and uncooperative institution designs: Surprises and problems in open-source game theory}
\author{Colomban Duclaux}
\date{\today}

\usepackage[backend=biber,style=alphabetic]{biblatex} % Ensure style matches your preference
\addbibresource{../references.bib} % Link your .bib file

\newcommand{\dupoc}{\texttt{DUPOC}}
\newcommand{\cupod}{\texttt{CUPOD}}
\newcommand{\cimcic}{\texttt{CIMCIC}}
\newcommand{\dimcid}{\texttt{DIMCID}}
\begin{document}

\maketitle
These notes recap the content exposed in \cite{critch2022cooperativeuncooperativeinstitutiondesigns}

\section{Abstract}
\textit{Open-source game theory}: game theory where the internal programming of each of the agents is visible.
\textit{Program equilibria}: game-theoretic equilibria in open-source game theory.

In this paper, give a series of counterintuitive results on open-source games.
Some of these open-source agents exhibit desirable characteristics: create incentives for cooperation.

\section{Introduction}
Important to understand how automated decision-making systems can interact.
Sometimes, counterintuitive strategic interactions between algorithms arise, that need to be checked-in.

Layout: \begin{itemize}
    \item Examine interactions between very simple agents who reason about each other using mathematical proofs and make decisions based on that reasoning.
    \item Show that results in mathematical logic can be applied to resolve the outcomes of such games.
    \item Conjecture that perhaps these results may be reducible to a form more easily applied to human institutions.
\end{itemize}

\section{Setup}
Agents will make decisions in part by formally verifying certain properties of each other: formal verifier agents. Examples: \texttt{CB}, \texttt{DB}, \texttt{CUPOD}, \texttt{DUPOC}

Proof-searching via proof-checking: \texttt{proof\_checker}, \texttt{proof\_search}


\begin{theorem} \label{thm:cupod-cupod}
    For $k$ large, $\texttt{outcome(CUPOD(k), CUPOD(k))} == (\texttt{D,D})$
\end{theorem}
\begin{proof}
    Proof is based on a modified version of Löb's Theorem. Essentially, it states that if (in our deduction system, we have that for all $k$, there exists a proof of above using less than $k$ characters implies above statement), then above statement.
\end{proof}

\subsection{\texttt{DUPOC}-like institutions}
Two \texttt{DUPOC}-like institutions that become open-source to one another can in turn cooperate together given the corollary of $\cref{thm:cupod-cupod}$.

Compared to \texttt{CUPOD}-like institutions that seem "nicer" in a sense because they never exploit its opponent, \texttt{DUPOC}-like institutions have desirable game-theoretic properties: \begin{enumerate}
    \item \textbf{unexploitability:} \dupoc never cooperates with an opponent that will defect
    \item \textbf{broad cooperativity:} \dupoc achieves (C, C) not just with other \dupoc s, but with other agents that follow "G-fairness" (opponents that always cooperates with agents that exhibit easy-to-verify cooperation)
\end{enumerate}

\section{Generalizability of results}
\begin{itemize}
    \item Non-deterministic version of \dupoc: \texttt{P}\dupoc which cooperates with probability $q$ if it knows its opponent cooperates with probability $q$.
    \item More efficient \texttt{proof\_search} methods
\end{itemize}

\section{More Results}
\cimcic: Cooperate if my cooperation implies cooperation from the opponent
$\to$ unexploitable.
\dimcid: defect if my cooperation implies defection from the opponent

\begin{theorem}
    \begin{itemize}
        \item For $k$ large, $\texttt{outcome}(\cimcic(k),\cimcic(k)) == (C,C)$
        \item For $k$ large, $\texttt{outcome}(\dupoc(k),\cimcic(k)) == (C,C)$
    \end{itemize}
\end{theorem}
\begin{theorem}
    \begin{itemize}
        \item For $k$ large, $\texttt{outcome}(\dimcid(k),\dimcid(k)) == (D,D)$
        \item For $k$ large, $\texttt{outcome}(\cupod(k),\dimcid(k)) == (D,D)$
    \end{itemize}
\end{theorem}

\section{Discussion}
\subsection{Applicability to non-logical reasoning processes}
Do agents need to use formal logic to form their beliefs, or can similar results be obtained with probabilistic reasoning?
Basically, we want to still have the PBLT so we need: \begin{itemize}
    \item  The proof language for representing the agents’ beliefs needs to be
expressive enough to represent numbers and computable functions and
to introduce and expand abbreviations.
    \item There must be a process the agents can follow for deriving beliefs from
other beliefs (writing a proof is such a process).
    \item The proof language must also be able to represent belief derivations
(e.g., proofs) as objects in some way, and refer to the cost (e.g., in time
or space) of producing various belief derivations.
\end{itemize}
Since these conditions are fairly general, it seems likely that similar results could be obtained for agents using probabilistic reasoning processes.
\subsection{Modelling self-fulfilling prophecies in humans}
\subsection{Comparison to Nash equilibria}

\section{Conclusion}
Open-source game theory is a rich field with many counterintuitive results, and should not be left unchecked.
One-shot interactions can "short-circuit" unbounded recursions of metacognition, like in \cref{thm:cupod-cupod}.
Still a bunch of open problems.



\printbibliography
\end{document}